% mode: ini
% WEB 1466, 1467: \uppercase and \lowercase.  The test t<cs_token_flag+single_base lets through characters of every category and
% active characters, and not control sequences.  equiv(b+c)=0 leaves a token alone.  The category stays; a token's character
% changes to any code, so an active character becomes another active character.  The list goes back by back_list: nested use,
% empty text, a text that defines something.  No expansion happens in the text.  Missing brace, and a runaway.
\input prelude
\errorcontextlines=0
\def\sh#1{\def\r{#1}\show\r}
\lccode`a=`a \lccode`b=`b \lccode`A=`a \lccode`B=`b
\uccode`a=`A \uccode`b=`B \uccode`A=`A \uccode`B=`B
\uppercase{\sh{abAB}}
\lowercase{\sh{abAB}}
\uppercase{\sh{xyz}}
\lowercase{\sh{XYZ}}
\uppercase{}
\lowercase{}
\uppercase{ }
\uppercase{\sh{}}
% control sequences stay, and their names
\uppercase{\sh{\a\ab\relax\undefinedone}}
\lowercase{\sh{\A\AB}}
\uppercase{\expandafter\sh\expandafter{\string\a}}
\def\m{a}\uppercase{\sh{\m}}
\uppercase\expandafter{\expandafter\sh\expandafter{\m}}
% a category stays
\catcode`\@=12
\uccode`\@=`a \uppercase{\sh{@}}\uccode`\@=0
\uccode`\$=`a \uppercase{\sh{$}}\uccode`\$=0
\uccode`\&=`a \uppercase{\sh{&}}\uccode`\&=0
\uccode`\#=`a \uppercase{\def\r{##}}\show\r \uccode`\#=0
\uccode`\^=`a \uppercase{\sh{^}}\uccode`\^=0
\uccode`\_=`a \uppercase{\sh{_}}\uccode`\_=0
\uccode`\ =`a \uppercase{\sh{ }}\uccode`\ =0
\uccode`\{=`a \uppercase{\sh{{}}}\uccode`\{=0
\uccode`\}=`a \uppercase{\sh{{}}}\uccode`\}=0
\uccode`\{=`x \uccode`\}=`y \uppercase{\sh{{a}}}\uccode`\{=0 \uccode`\}=0
\uccode`1=`9 \uppercase{\sh{1}}\uccode`1=0
\uccode`a=`1 \uppercase{\count1=a \showthe\count1 \sh{a}}\uccode`a=`A
\uccode`a=0 \uppercase{\sh{a}}\uccode`a=`A
\lccode`A=0 \lowercase{\sh{A}}\lccode`A=`a
\uccode`a=255 \uppercase{\sh{a}}\uccode`a=`A
% an active character
\catcode`\!=13 \def!{bang}
\def~{tilde}
\uccode`~=`! \uppercase{\sh{~}}
\uppercase{\show~}
\uccode`~=0 \uppercase{\sh{~}}
\lccode`!=`~ \lowercase{\show!}\lowercase{\sh{!}}
\lccode`!=0
\catcode`\!=12
\catcode`\^^A=13 \uccode`\^^A=`\^^B \uppercase{\sh{^^A}}\uppercase{\show^^A}\uccode`\^^A=0
\catcode`\^^A=15
% a token that takes effect
\uppercase{\let\z=a}\show\z
\lowercase{\let\z=A}\show\z
\uppercase{\def\z{abc}}\show\z
\lowercase{\global\def\z{ABC}}\show\z
\uppercase{\uppercase{\sh{a}}}
\uppercase{\lowercase{\sh{a}}}
\lowercase{\uppercase{\sh{a}}}
\uppercase{\begingroup}\endgroup
\uppercase{\count1=1\relax}\showthe\count1
\uccode`0=`5 \uppercase{\count1=10\relax}\showthe\count1 \uccode`0=0
\uppercase{\show a}
\uppercase{\show\relax}
% an \edef does not run it; \the before it does
\edef\z{\uppercase{a}}\show\z
\toks0={a}\uppercase\expandafter{\the\toks0}\par
\edef\z{\noexpand\uppercase{\the\toks0}}\show\z
% what scan_toks does: no expansion, a brace needed, implicit braces, a macro in the way
\def\t{abc}\uppercase{\t}\par
\let\bg={\uppercase\bg \sh{a}}
\uppercase\relax{\sh{b}}
\uppercase x}\par
\uppercase\t}\par
\uppercase}\par
\long\uppercase{a}\par
\global\lowercase{a}\par
\outer\def\o{}
\uppercase{\o}\par
\def\o{}
\uppercase{\par}
\uppercase{a\par b}\par
\showlists
\uppercase{ this one never ends
\end
